% ============================================
% Methods
% ============================================
\section{Methods}
\label{sec:methods}

% --------------------------------------------
\subsection*{Microbial Strain Library, Media, and Culture}
\label{sec:methods:strains}]

We obtained a library of 54 bacterial isolates derived from soil, tree leaves, and flower environments collected in Cambridge, MA, USA. Each isolate was purified by serial streaking and stored as glycerol stocks (25\% glycerol, $-80$~°C). Phylogenetic identities were determined by 16S rRNA gene sequencing of the V4 region and BLAST alignment to the SILVA 138 database; the 54 isolates corresponded to 50 unique amplicon sequence variants (ASVs) spanning 29 families across three phyla (Proteobacteria, Firmicutes, and Bacteroidota) (Supplementary Fig.~S1). 

The culture medium (CM) for isolates was composed of 1~g~L$^{-1}$ yeast extract and 1~g~L$^{-1}$ soytone (both from Becton Dickinson), 10~mM sodium phosphate buffer adjusted to pH~6.5, 0.1~mM CaCl$_2$, 2~mM MgCl$_2$, 4~mg~L$^{-1}$ NiSO$_4$, and 50~mg~L$^{-1}$ MnCl$_2$. All media were filter sterilized using bottle-top filtration units (VWR), and all chemicals were purchased from Sigma--Aldrich unless otherwise noted. 

Our coalescence experiments were conducted in three media conditions with varying nutrient supplementation, which produce different strengths of pairwise interactions: (1) \textbf{Nutr$-$}, the culture medium without additional carbon or nitrogen supplementation; (2) \textbf{Base}, CM supplemented with 5~g~L$^{-1}$ glucose and 4~g~L$^{-1}$ urea; and (3) \textbf{Nutr+}, CM supplemented with 20~g~L$^{-1}$ glucose and 16~g~L$^{-1}$ urea. Both parental community assembly and coalescence experiments were performed in the same media condition throughout each experimental series.

Both monocultures and communities were grown in 300~$\mu$L volumes in 96-well deep-well plates (Deepwell plate 96/1000~$\mu$L; Eppendorf) covered with AeraSeal adhesive sealing films (Excel Scientific). Plates were incubated aerobically at 25~°C and shaken at 800~rpm on Titramax shakers (Heidolph). To minimize pre-adaptation, all isolates from frozen stocks were preconditioned for two serial growth--dilution cycles prior to assembling parental communities. 
%Each isolate's growth rate and final optical density were measured in monoculture, confirming broad variation in growth yield and pH modification potential (Supplementary Fig.~S30 and Supplementary Methods).

% --------------------------------------------
\subsection*{Construction of Parental Communities and Community Coalescence Experiments}
\label{sec:methods:coalescence}

Parental communities were constructed by sequential assignment of isolates from the strain library at three richness levels: 6, 12, or 24 members (pool sizes P6, P12, and P24). In total, 42 distinct parental communities were assembled (9 $\times$ P6, 9 $\times$ P12, 12 $\times$ P24), each with two biological replicates. For P6 and P12 communities, 9 non-overlapping communities were assembled by sequential assignment (e.g., community 1: strains 1--6, community 2: strains 7--12, etc.). For P24 communities, because only two non-overlapping sets of 24 species can be drawn from a pool of 54 isolates, we assembled 12 partially overlapping communities to enable sufficient coalescence pairing diversity. Parental communities were inoculated by mixing equal volumes of stationary-phase monocultures of each constituent isolate, then stabilized by seven daily serial dilutions ($\times$30 every 24~h) prior to coalescence assays.

Coalescence experiments were conducted by mixing two pre-stabilized parental communities (A and B) at equal volume ratio (1:1) after seven daily dilution cycles. Coalescence pairs were composed systematically to maximize coverage of community combinations while maintaining experimental feasibility. For P6 and P12 richness levels, we performed near-complete pairwise coalescence: P6 included 14 coalescence pairs from 9 parental communities, and P12 included 27 coalescence pairs from 9 parental communities. For P24, due to the larger number of possible combinations, we performed a subset of 6 coalescence pairs using non-overlapping parental community pairs, resulting in 47 total coalescence pairs for synthetic communities. Experiments were conducted in all three media conditions (Nutr$-$, Base, and Nutr+), and each condition was analyzed separately; some coalescence events were excluded due to sequencing failures. The resulting mixtures were cultured for another seven serial transfers ($\times$30 dilutions every 24~h) under identical media conditions to reach new steady states (\figref{fig:fig1}B).


% --------------------------------------------
\subsection*{Measuring Community Composition via 16S rRNA Sequencing}
\label{sec:methods:sequencing}

To measure the composition of parental communities and coalesced communities, we used 16S ribosomal RNA (rRNA) amplicon sequencing. DNA extraction was performed with the QIAGEN DNeasy PowerSoil HTP 96 Kit following the protocol provided by the manufacturer. The obtained DNA was used for 16S (V4 region) amplicon sequencing. Library preparation and Illumina MiSeq sequencing were performed by the Environmental Sample Preparation and Sequencing Facility at Argonne National Laboratory. We used DADA2\citep{Callahan2016} to obtain amplicon sequence variants (ASVs). Taxonomic identities were assigned to ASVs using SILVA (version 138)\citep{Quast2013} as a reference database. For each sample, species richness was calculated as the number of ASVs with a relative abundance $\geq$0.1\%, which corresponds to the 0.1\% extinction threshold used in simulations. The phylogenetic tree (Supplementary Fig.~S2) was constructed using Simple Phylogeny\citep{Madeira2022} by the EMBL's European Bioinformatics Institute. All plots of relative ASV abundances with stacked bars show results from one representative replicate. All raw reads are available in the NCBI Sequence Read Archive under BioProject accession PRJNAXXXXXX.
% \jy{NOTE: Technical replicate correlation was previously referenced as Fig S4, but S4 is similarity metric robustness. Need to add a separate figure for technical replicates or remove reference.}


% --------------------------------------------
\subsection*{Similarity-Based Classification of Coalescence Outcomes}
\label{sec:methods:classification}

To characterize coalescence outcomes, we quantified the compositional similarity between the post-coalescence community and each parental community. For each coalescence event, let $\vec{x}_A$, $\vec{x}_B$, and $\vec{x}_C$ represent the L$_2$-normalized abundance vectors of parental communities A and B and the coalesced community C, respectively. We computed the cosine similarity between the coalesced community and each parent:
\begin{equation}
\text{Sim}(C, A) = \vec{x}_C \cdot \vec{x}_A, \quad \text{Sim}(C, B) = \vec{x}_C \cdot \vec{x}_B
\label{eq:similarity_methods}
\end{equation}
These two similarity scores place each coalescence outcome in a two-dimensional similarity space, where the position reflects how strongly the coalesced community resembles each parent (\figref{fig:fig1}B).

Because parental communities may share species (overlapping ASVs), the parental vectors $\vec{x}_A$ and $\vec{x}_B$ are generally not orthogonal ($\vec{x}_A \cdot \vec{x}_B \neq 0$). To account for this non-orthogonality, we computed the contribution coefficients $(u, v)$ by projecting $\vec{x}_C$ onto the subspace spanned by $\vec{x}_A$ and $\vec{x}_B$:
\begin{equation}
\begin{pmatrix} u \\ v \end{pmatrix} = \begin{pmatrix} \vec{x}_A \cdot \vec{x}_A & \vec{x}_A \cdot \vec{x}_B \\ \vec{x}_A \cdot \vec{x}_B & \vec{x}_B \cdot \vec{x}_B \end{pmatrix}^{-1} \begin{pmatrix} \vec{x}_C \cdot \vec{x}_A \\ \vec{x}_C \cdot \vec{x}_B \end{pmatrix}
\end{equation}
When parents share no species ($\vec{x}_A \cdot \vec{x}_B = 0$), the coefficients reduce to the cosine similarities: $u = \text{Sim}(C, A)$ and $v = \text{Sim}(C, B)$. Negative coefficients were set to zero since negative contributions are biologically meaningless. The residual magnitude $k = \|\vec{x}_C - u \cdot \vec{x}_A - v \cdot \vec{x}_B\|$ quantifies the component of the coalesced community orthogonal to both parents, representing novel restructuring.

From the coefficients $(u, v, k)$, we derived two classification metrics:

Retention magnitude:
\begin{equation}
x = \sqrt{u^2 + v^2}
\end{equation}
This metric quantifies the fraction of the coalesced community composition explained by the parental communities. Values of $x$ close to 1 indicate high retention of parental composition; values close to 0 indicate extensive restructuring.

Parental Dominance Index (PDI):
\begin{equation}
\text{PDI} = \frac{\arctan(v/u) - \pi/4}{\pi/4}
\end{equation}
This metric quantifies the selection preference toward one parent over the other. When $u = v$ (equal contributions), PDI $= 0$. PDI ranges from $-1$ (complete dominance by parent A, when $v = 0$) to $+1$ (complete dominance by parent B, when $u = 0$).

Based on these metrics, each coalescence event was classified into one of three outcome categories:

\begin{itemize}
    \item Restructuring: $x^2 \leq 0.5$. Less than half of the coalesced community composition is explained by the parental communities, indicating substantial ecological reorganization.
    \item Mixture: $x^2 > 0.5$ and $|\text{PDI}| \leq 0.5$. High retention with balanced contributions from both parents, indicating coexistence of parental lineages.
    \item CLS (community-level selection): $x^2 > 0.5$ and $|\text{PDI}| > 0.5$. High retention but one parent dominates, indicating selection for one parental community over the other.
\end{itemize}

For sensitivity analyses of classification boundaries and alternative similarity metrics, see Supplementary Information.


% --------------------------------------------
\subsection*{Lotka--Volterra Simulations}
\label{sec:methods:simulations}

We modeled species-abundance dynamics with the generalized Lotka--Volterra (gLV) system:
\begin{equation}
\frac{dN_i}{dt} = g_i N_i \left(1 - \frac{1}{k_i} \sum_j I_{ij} N_j \right)
\label{eq:glv_methods}
\end{equation}
where $N_i(t)$ is the abundance of species $i$ at time $t$, $g_i$ is its intrinsic growth rate, $k_i$ is its carrying capacity, and $I_{ij}$ is the pairwise interspecies interaction coefficient between species $i$ and $j$. To isolate the role of interaction strength, we fixed $g_i = 1$ and $k_i = 1$ for all species in all simulations. The interaction matrix $\mathbf{I}$ had $I_{ii} = 1$ (intraspecific competition) and $I_{ij} \sim \text{Unif}(0, 2\mu)$ for $i \neq j$, with $\mu$ denoting the mean interspecific interaction strength.

Each simulation used a pool of $N = 48$ species partitioned into four non-overlapping ``parental'' communities of 12 species. For each replicate, we randomly permuted the 48 species and assigned them sequentially (1--12, 13--24, 25--36, 37--48) to form the four communities, ensuring no species were shared across parental communities, analogous to the P6 and P12 experimental design. A fresh interaction matrix $\mathbf{I}$ was independently sampled for every replicate so that coalescence dynamics were explored across diverse ecological contexts rather than a single fixed species pool.

For single-community assembly, we numerically integrated the gLV equations from random initial conditions $N_i(0) \sim \text{Unif}(0, 0.1)$ to ecological steady state. We solved the ODEs with the adaptive Runge--Kutta method RK23 implemented in \texttt{scipy.integrate.solve\_ivp}, integrating from $t = 0$ to $t = 5000$, which was sufficient for all communities to equilibrate. Following integration, species with abundances below an extinction threshold of $10^{-4}$ were set to zero, yielding four parental-community steady states per replicate.

Community coalescence was simulated by mixing pre-equilibrated parental communities at equal proportions, mirroring the experimental 1:1 volume protocol. For each of the $\binom{4}{2} = 6$ parental community pairs, we retrieved the steady-state abundance vectors $\mathbf{y}^{(1)}$ and $\mathbf{y}^{(2)}$, formed the equal mixture $\mathbf{y}_{\text{mix}} = (\mathbf{y}^{(1)} + \mathbf{y}^{(2)})/2$, and zeroed species falling below the extinction threshold in this initial mixture. We then re-integrated the gLV system from $\mathbf{y}_{\text{mix}}$ using the same interaction matrix $\mathbf{I}$ over $t \in [0, 5000]$ until a new steady state was reached, and reapplied the $10^{-4}$ threshold to the final abundances. This procedure captures the ecological relaxation after mixing and allows competitive exclusion, coexistence, or the emergence of previously rare species to arise from the underlying interaction structure.


% --------------------------------------------
\subsection*{Pairwise Invasion Assays}
\label{sec:methods:invasion}

To empirically estimate interaction strength, we performed pairwise invasion experiments among the 12 most abundant isolates. Isolate pairs were mixed at a 95:5 ratio (resident:invader) by colony count and propagated in Nutr$-$, Base, and Nutr+ media under daily dilution ($\times$30) for 4 cycles. Final compositions were determined by colony counting on agar plates. Competitive outcomes were scored as coexistence, exclusion, or bistability based on final relative abundances (coexistence if both $>$10\%; exclusion if one $<$1\%). The fraction of exclusions was used as a proxy for mean negative interaction strength under each nutrient condition (\figref{fig:fig4}B).


% --------------------------------------------
\subsection*{Statistical Analyses}
\label{sec:methods:statistics}

All analyses were performed in Python 3.11 using NumPy, pandas, and scikit-learn, with statistical significance defined at $\alpha = 0.05$ unless otherwise noted.

\begin{itemize}
    \item Paired $t$-tests were used to compare mean interaction strengths before and after community assembly (\figref{fig:fig2}C).
    \item Permutation tests (1,000 permutations) were used to assess whether pairwise selection correlations differed between same-parent and cross-parent species pairs (\figref{fig:fig2}D, Supplementary Fig.~S21). The null distribution was generated by shuffling species origin labels within each coalescence event.
    \item Mann-Whitney U tests were used to compare experimental one-sided selection values against null model distributions (Supplementary Fig.~S8).
    \item Linear regression was used to quantify the predictability of coalescence outcomes from dominant-species pairwise competition, with $R^2$ reported as the coefficient of determination (\figref{fig:fig5}C).
    \item $\chi^2$ tests were used to compare observed vs.\ expected counts of CLS, Mixture, and Restructuring events.
    \item Mixed-effects models (lme4 package in R) were used to control for parental community richness and community identity effects when quantifying nutrient dependence.
\end{itemize}

For detailed descriptions of skewness null model analysis, sensitivity analysis, and predictability analysis, see Supplementary Information.


% --------------------------------------------
\subsection*{Obtaining Natural Sample-Derived Communities and Coalescence}
\label{sec:methods:natural}

To test whether the coalescence patterns observed in synthetic communities generalize to more ecologically relevant settings, we performed additional coalescence experiments using communities derived from natural environmental samples. Unlike the synthetic communities above, which were assembled from individually isolated strains, natural sample-derived communities retain the complex species assemblages and ecological interactions present in their native habitats. Six environmental samples were collected from diverse microhabitats in Cambridge, MA, USA, including soil, compost, and decomposing organic matter. Each sample was homogenized by vortexing and inoculated into culture medium (CM), then subjected to seven serial growth-dilution cycles (24~h per cycle, $\times$30 dilution) to enrich for laboratory-cultivable bacteria and establish stable communities prior to coalescence.

Coalescence experiments followed the same protocol as synthetic communities. After stabilization, 16 pairwise coalescence events were performed from the 6 parental communities, mixed at 1:1 volume ratio and allowed to re-equilibrate through seven additional serial growth-dilution cycles. Experiments were conducted in all three media conditions (Nutr$-$, Base, and Nutr+) to test whether interaction strength modulates coalescence outcomes in natural sample-derived communities as in synthetic communities. Community composition at each stage (parental and post-coalescence) was characterized via 16S rRNA amplicon sequencing using the same protocol as synthetic communities (Supplementary Fig.~S28).
